
\documentclass{article}
\usepackage[style=apa]{biblatex}
\DeclareLanguageMapping{american}{american-apa}
\addbibresource{synthetic_other_apa7_clean.bib}

\title{A Synthetic Manuscript for Reference-Extraction Benchmarking}
\author{Synthetic Author}
\date{2026}

\begin{document}
\maketitle

\begin{abstract}
This synthetic manuscript exists only to exercise the reference-extraction pipeline against a document with realistic structure and an exactly known bibliography, generated from the entries in entries.py.
\end{abstract}

\section{Introduction}
The question taken up here has been approached from several directions, and the paragraphs below trace the line of work this manuscript builds on rather than restating each contribution in full. A closely comparable argument is developed by \textcite{higgs1964broken}.

\section{Method}
The procedure follows established practice at every step, so what matters in this section is which earlier report each decision was taken from.

\section{Results}
The pattern observed in these data is reported first and then placed against the findings it is meant to be read alongside.

\section{Discussion}
What remains is to say which of the earlier claims the present data support and which of them they leave open.

\printbibliography

\end{document}
